% ============================================================================
%  Chapter 0 -- Course Introduction                              (Lecture 00)
% ============================================================================
\chapter{Introduction: Program First, Formalise Later}
\label{ch:intro}

\epigraphline{The Matrix is a system, Neo. That system is our enemy.}
  {\emph{The Matrix}, Lana and Lilly Wachowski, 1999}

\begin{chapgoals}
\begin{itemize}
\item What the course is, who teaches it and how its ninety-six hours are used.
\item The method: a program that runs comes first, its theory comes after.
\item The machine: \prg, \riscv and the debugger \prgqt.
\item The project: \vermone, a game written from an empty file and published.
\item The rules: syllabus, learning outcomes, assessment, material.
\end{itemize}
\end{chapgoals}

A computer is the most complicated object that most people use every day, and
the one whose inner working they are least encouraged to look at. Its makers
have spent seventy years hiding the machine behind layers that are each easier
to use than the one below: logic gates behind instructions, instructions behind
languages, languages behind libraries, libraries behind applications, and now
applications behind conversational assistants. The layers are what make
computing useful. They are also a system in the sense of the film quoted
above: one can live inside it for a lifetime without knowing that there is
anything underneath.

Computer architecture is the discipline that looks underneath. It studies how
a machine that can only switch currents on and off comes to represent numbers,
text, pictures and sound; how a handful of logical operations, repeated
thousands of millions of times a second, become arithmetic, decisions and
memory; and how the whole is organised so that a sequence of very simple
orders, a \term{program}, can make it do anything that can be computed at all.
This chapter explains how the course approaches that study, with which machine
and through which project, and then states the practical rules.

% ----------------------------------------------------------------------------
\section{The course at a glance}

\emph{Computer Architecture and Computer Architecture Laboratory}, code ARC12,
is a first-year, first-semester course of the Bachelor's Degree in Computer
Science of the University of Naples ``Parthenope''. It is worth 12~ECTS
credits, belongs to the scientific area INFO-01/A (Informatics) and consists
of 96~hours, divided equally between lectures and laboratory
(Table~\ref{tab:intro:glance}).

\begin{table}[htb]
\centering\small
\caption{The course at a glance.}\label{tab:intro:glance}
\begin{tabularx}{\textwidth}{@{}lL@{}}
\toprule
Name & Computer Architecture and Computer Architecture Laboratory (ARC12)\\
Position & First year, first semester, Bachelor's Degree in Computer Science\\
Credits & 12 ECTS, 96 hours; scientific area INFO-01/A (Informatics)\\
Lectures & Raffaele Montella, 6 ECTS, 48 hours\\
Laboratory & Mariacarla Staffa, 6 ECTS, 48 hours\\
Office hours & Room 420, Tuesday 14:00--16:00, by appointment; also on
  Microsoft Teams (team code \code{c7rf33g})\\
Extra credits & Accenture SuperSkills! gives 2 ECTS as further knowledge
  (\url{https://richiediaccesso.uniparthenope.it})\\
Prerequisites & None\\
\bottomrule
\end{tabularx}
\end{table}

There are no formal prerequisites. No previous experience of programming, of
electronics or of mathematics beyond the school curriculum is assumed: the
only real prerequisite is a passion for computing and for the technologies that
make it possible.

The two halves of the course are two views of the same matter. In the
\textbf{lectures} the ideas are presented: information, numbers, Boolean
algebra, logic networks, memory, the processor, input and output. Every topic
opens with a program that can be run and is then formalised with definitions,
methods and proofs; every lecture closes with multiple-choice questions in the
style of the midterm test; and every four lectures a tutorial adds the next
part of \vermone, the game of the course. In the \textbf{laboratory} the same
ideas are practised: number representation and conversion techniques;
analysis, synthesis and minimisation of logic circuits; \riscv assembly
programming on the \prg platform; and, finally, the student's own project, a
variant of \vermone or a game of the student's invention. This book is the
text of the lectures; its exercises are the natural material of the
laboratory.

% ----------------------------------------------------------------------------
\section{A new approach}

\begin{prgconcept}[Program first. Formalise later.]
From the first days the student writes \riscv assembly that runs on a real
processor. Each formal topic then answers a question that the student's own
program has already raised.
\end{prgconcept}

The traditional order of the subject is deductive: representation, then logic,
then circuits, then the processor, then, at last, a program. It is the order
in which the machine is \emph{built}, and it has the elegance of every
construction from first principles. It is not the order in which the machine
is \emph{understood}. A student who meets two's complement in the third week
of a deductive course has never seen a negative number inside a computer, and
has no reason to care how it is written. A student who, the week before, has
moved a rectangle to the left of the screen until its coordinate, shown by the
debugger, has turned from \hex{00000000} into \hex{FFFFFFFF} has a question,
and two's complement is its answer.

The course is therefore organised as a spiral, repeated for every topic
(Figure~\ref{fig:intro:spiral}).

\begin{figure}[htb]
\centering
\begin{tikzpicture}[node distance=7mm]
  \node[blkacc,text width=27mm,minimum height=20mm] (a) {\textbf{Play}\\[2pt]\footnotesize type a few instructions, build a cartridge, run it};
  \node[blkamb,text width=27mm,minimum height=20mm,right=of a] (b) {\textbf{Observe}\\[2pt]\footnotesize a number wraps, a bit lights a lamp, a rectangle moves};
  \node[blkvio,text width=27mm,minimum height=20mm,right=of b] (c) {\textbf{Formalise}\\[2pt]\footnotesize definitions, theorems and methods explain why};
  \node[blkgrn,text width=27mm,minimum height=20mm,right=of c] (d) {\textbf{Return}\\[2pt]\footnotesize predict the result before running it};
  \draw[->,thick,prgink] (a) -- (b); \draw[->,thick,prgink] (b) -- (c); \draw[->,thick,prgink] (c) -- (d);
  \draw[->,thick,prgink,rounded corners=3mm] (d.south) -- ++(0,-7mm) -| node[pos=0.25,below,lbl]{next topic} (a.south);
\end{tikzpicture}
\caption{The learning spiral, repeated for every topic of the course.}
\label{fig:intro:spiral}
\end{figure}

First one \emph{plays}: a few instructions are typed, a cartridge is built and
run. Then one \emph{observes}: something happens that was not expected, or was
expected and can now be seen. Then one \emph{formalises}: definitions and
theorems explain exactly why. Finally one \emph{returns} to the machine,
and the measure of understanding is the ability to predict a result before
running the program. A mark is an opinion; a paddle that leaves the screen is
a fact. The course uses facts, and then explains them.

\subsection{What the course aims at}

The aims are those of every serious course on the subject; only the path is
different. The course

\begin{itemize}
\item presents the fundamental aspects of the architecture of modern computers;
\item uses, in the laboratory, the basic concepts needed to program them
  efficiently in low-level languages: machine language and assembly;
\item treats rigorously the tools from which a processor is built, namely
  Boolean algebra with the synthesis of combinational and sequential logic
  networks, and binary representation with binary arithmetic;
\item describes in detail the functional and physical structure of the CPU, of
  the memory and of the input/output subsystem;
\item develops assembly code that makes visible how a processor really works
  and how the basic units of a computer interact.
\end{itemize}

% ----------------------------------------------------------------------------
\section{The machine: PRG32 and RISC-V}

\subsection{A retro console programmed in assembly}

\prg is an open-source educational runtime for teaching \riscv assembly and C,
developed at the University of Naples ``Parthenope'' \citep{prg32repo}. To
the eye it is a small games console of the kind that was sold in the early
1980s (Figure~\ref{fig:intro:prg32}); to the programmer it is a deliberately
thin contract.

\begin{figure}[htb]
\centering
\begin{subfigure}[b]{0.48\textwidth}\centering
\includegraphics[width=\textwidth]{menu_main.png}
\caption{The menu of the resident firmware.}
\end{subfigure}\hfill
\begin{subfigure}[b]{0.48\textwidth}\centering
\includegraphics[width=\textwidth]{game_pong.png}
\caption{Pong: native \riscv machine code.}
\end{subfigure}
\caption{\prg, the machine of the course.}
\label{fig:intro:prg32}
\end{figure}

\begin{itemize}
\item The student's code runs \emph{natively}, as 32-bit \riscv machine code
  of the RV32IMAC instruction set. It is not interpreted by a teaching
  simulator: the instructions written are the instructions executed.
\item The contract consists of three functions that the student writes:
  \code{init} runs once; \code{update} and \code{draw} run once per frame.
\item Games travel as small \term{cartridges}, files with the extension
  \fname{.prg32}.
\item The machine offers a viewport of 320 by 200 pixels, a joystick with a
  few buttons, and sound: small enough to be held in one's head, rich enough
  to exercise the whole course.
\end{itemize}

\subsection{The philosophy behind PRG32}

Six principles, taken from the textbook of the platform
\citep{prg32textbook}, shape both \prg and the way this book uses it.

\begin{description}[style=nextline,leftmargin=1.5em]
\item[Small enough to understand.] One processor, a screen of 320 by 200
  pixels, a joystick, a small speaker. Nothing is hidden: the whole machine can be
  held in one's head.
\item[Real, not simulated.] The instructions written are those that the chip
  executes. What is learnt here is true on every \riscv processor, from a
  sensor to a supercomputer.
\item[Limits make you creative.] A few kilobytes and three functions. The
  games of the 1980s were born from limits like these: a constraint is an
  invitation.
\item[One thing at a time.] Start with functions that only return. Add one
  thing, run, look. \code{update} changes the state, \code{draw} shows it.
  Sound comes last.
\item[Look inside.] When something is surprising, stop the program and read
  registers and memory. The debugger is part of the method, not an emergency
  tool.
\item[Make, then share.] A cartridge is a file that can be given to a friend.
  The Cartridge Store is the shelf where finished work meets its players.
\end{description}

The platform has, accordingly, two uses in the course. It is a
\emph{gymnasium}: every chapter opens with a short program that runs and
closes by going back to the machine; an exercise consists of a prediction and
of a check; and the machine never flatters, because a wrong mask lights the
wrong lamp. And it is a \emph{workshop}: a blank cartridge and a screen, on
which rules, colours, sounds and levels are design decisions that belong to
the author. The project of the course starts guided and ends free, and what
is finished can be published. First-year students rarely expect a course on
architecture to ask them for an idea; this one does.

\subsection{PRG32-QT: the debugger of the course}

\prgqt is \prg on a laptop: an application that runs the same cartridges with
no board and no compiler \citep{prg32qt}. In debug mode it opens the machine
(Figure~\ref{fig:intro:debugger}): the program can be paused, resumed and
advanced one instruction at a time; breakpoints can be placed on any
instruction of \code{init}, \code{update} or \code{draw}; the thirty-two
registers and 128~bytes of memory at any address are displayed; the speed can
be set between 0.01 and 4 times the real one; and counters report frames,
instructions executed and frames that arrived late.

\begin{figure}[htb]
\centering
\includegraphics[width=0.92\textwidth]{prg32qt_debugger.png}
\caption{\prgqt in debug mode: the game on the left, the machine on the right.}
\label{fig:intro:debugger}
\end{figure}

The debugger is used in every chapter: to check a prediction, to find a
mistake, to measure. It also supplies the working definition of understanding
adopted in this book.

\begin{prgconcept}[When is a program understood?]
A program is understood when one can say what each register will hold
\emph{before} pressing \key{Step}.
\end{prgconcept}

\subsection{Why RISC-V?}

The choice of the processor is not neutral, and four reasons led to \riscv.

\begin{description}[style=nextline,leftmargin=1.5em]
\item[It is open.] \riscv is an open standard instruction set: anyone may
  design, build and sell a \riscv processor without paying a licence for the
  architecture \citep{riscvisa}.
\item[It is small.] The base integer instruction set has about forty
  instructions and thirty-two registers. All of it can be learnt in one
  semester, and Appendix~\ref{app:riscv} lists all of it.
\item[It is real.] \ins{addi}, \ins{lw}, \ins{sw} and \ins{beqz} are not a
  teaching dialect. Industry ships \riscv cores by the billion, and they
  execute exactly these instructions.
\item[It is European.] Europe treats \riscv as a lever for technological
  sovereignty: the European Chips Act and projects such as DARE (Digital
  Autonomy with RISC-V in Europe, a EuroHPC Joint Undertaking project) build
  processors on it \citep{eurohpcdare}.
\end{description}

\subsection{The map: every topic is met twice}

Table~\ref{tab:intro:map} is the course on one page. Each formal topic of the
left column is met first on the machine, in the form given by the right
column, and only then in its general form. It is worth coming back to this
table at the beginning of each part of the book.

\begin{table}[htb]
\centering\small
\caption{Every formal topic is met first on \prg.}\label{tab:intro:map}
\begin{tabularx}{\textwidth}{@{}P{0.36\textwidth}L@{}}
\toprule
\textbf{Formal topic} & \textbf{First met on \prg as\ldots}\\
\midrule
Positional notation, hexadecimal & the colour \hex{07E0} and the value shown by the debugger\\
Two's complement, overflow & a position that moves left past zero and shows \hex{FFFFFFFF}\\
Fixed and floating point & a ball that moves 0.375 pixels per frame on a chip with no floating-point unit\\
Boolean algebra, bit masks & \code{andi t3, t0, 1} to test whether \key{LEFT} is pressed\\
Electronics, logic gates & the button that joins a pin of the chip to ground\\
Combinational networks, adders & what \ins{add} actually does to thirty-two bits\\
Sequential networks, registers & the counter that survives from one frame to the next\\
CPU, datapath, instruction cycle & single-stepping one instruction in \prgqt\\
Memory and input/output & \ins{lw} and \ins{sw}, the stack, the joystick and the display\\
\bottomrule
\end{tabularx}
\end{table}

% ----------------------------------------------------------------------------
\section{The project: Verm1}

\vermone, read \emph{vermone}, ``the big worm'', is a tribute to Snake: the
game of the arcade machine \emph{Blockade} (1976) and of the Nokia telephones
(1997). A worm moves on a grid, eats, grows, and must not bite itself. The
student writes every line of it in \riscv assembly, with no engine and no
library beyond the calls of \prg. The game is small enough for a first-year
student and rich enough to need every topic of the course; at the end it is
ready for the Cartridge Store with its author's name on it. The reference
implementation, against which every reader can compare a program of their
own, was contributed by Giuseppe Coviello (Appendix~\ref{app:verm1}).

The project is the thread of the book. It is built in four parts and a
publication (Table~\ref{tab:intro:verm1}, Figure~\ref{fig:intro:verm1}), and
each part uses exactly what the four chapters before it have taught. Each
part ends with a game that runs: the reader never goes home with a broken
program.

\begin{table}[htb]
\centering\small
\caption{\vermone: four parts and a publication.}\label{tab:intro:verm1}
\begin{tabularx}{\textwidth}{@{}lP{0.20\textwidth}LP{0.22\textwidth}@{}}
\toprule
\textbf{Chapter} & \textbf{After the part on} & \textbf{\vermone gains} & \textbf{With}\\
\midrule
\ref{ch:verm1p1}\quad Part 1 & information and numbers & a grid, a head that moves, a world without edges & colours, shifts, fixed point\\
\ref{ch:verm1p2}\quad Part 2 & logic and circuits & steering, food, a score & masks, XOR, tables\\
\ref{ch:verm1p3}\quad Part 3 & state and memory & a body that grows, collisions, three screens & a state machine, arrays\\
\ref{ch:verm1p4}\quad Part 4 & the machine & digits, levels, sound, fast drawing & calls, division, the frame buffer\\
\ref{ch:store}\quad The Store & the whole course & a name, an icon, credits, a bundle & the \prg tools\\
\bottomrule
\end{tabularx}
\end{table}

\begin{figure}[htb]
\centering
\begin{subfigure}[b]{0.24\textwidth}\includegraphics[width=\textwidth]{verm1_p1.png}\caption{Part 1: a head}\end{subfigure}\hfill
\begin{subfigure}[b]{0.24\textwidth}\includegraphics[width=\textwidth]{verm1_p2.png}\caption{Part 2: food, score}\end{subfigure}\hfill
\begin{subfigure}[b]{0.24\textwidth}\includegraphics[width=\textwidth]{verm1_p3_play.png}\caption{Part 3: a body}\end{subfigure}\hfill
\begin{subfigure}[b]{0.24\textwidth}\includegraphics[width=\textwidth]{verm1_p4_play.png}\caption{Part 4: finished}\end{subfigure}
\caption{\vermone grows with the reader: the same file, four sittings. The
screenshots are taken from the reference programs by the automatic tests of
the course material.}
\label{fig:intro:verm1}
\end{figure}

A tutorial chapter obeys four rules. It is \emph{self-contained}: it shows the
whole program, and every line is either new, changed, or reprinted unchanged.
It proceeds \emph{step by step}: a few lines at a time, and after each step
the program is built, run and compared with a stated result. It uses
\emph{the debugger}: each tutorial has a session in \prgqt, with a breakpoint,
a register to watch and a region of memory to read. And its result belongs to
the reader: each tutorial ends with changes to be made independently, because
the game that is published should not be identical to the one printed here.

% ----------------------------------------------------------------------------
\section{The course in blocks}

The order of the topics was chosen for learning. First comes what a bit
means; then how bits are combined; then how they are remembered; then the
machine that does all three. The processor is presented before the devices
that it drives. After each block the project uses it. The last lecture is an
outlook, held before a visit to the computing systems of the Department
(Table~\ref{tab:intro:blocks}).

\begin{table}[htb]
\centering\small
\caption{The blocks of the course. Numbers are those of the lectures and of
the chapters of this book.}\label{tab:intro:blocks}
\begin{tabularx}{\textwidth}{@{}lLP{0.28\textwidth}@{}}
\toprule
\textbf{Block} & \textbf{Lectures} & \textbf{Then}\\
\midrule
Start & 0 Introduction; 1 \prg and \prgqt: setting up; 2 First contact with \riscv & \\
Information & 3 History; 4 Digital information; 5 Integers; 6 Reals & 7 \vermone, part 1\\
Logic & 8 and 9 Boolean algebra; 10 Electronics; 11 Combinational machines & 12 \vermone, part 2\\
State and memory & 13 Sequential machines; 14 Registers and buses; 15 Memory; 16 Permanent memory & 17 \vermone, part 3\\
The machine & 18 Processor; 19 Instructions and interrupts; 20 Video card; 21 Input and output & 22 \vermone, part 4\\
Publication & & 23 \vermone on the Cartridge Store\\
Outlook & 24 Beyond one core: superscalar \riscv, HPC, GPUs and FPGAs & the visit to the computing systems\\
\bottomrule
\end{tabularx}
\end{table}

% ----------------------------------------------------------------------------
\section{Syllabus and learning outcomes}

\subsection{Syllabus}

\begin{description}[style=nextline,leftmargin=1.5em]
\item[Computer organisation.] Hardware and software; application and system
  software; the von Neumann model; operating principles of a computer.
\item[Boolean algebra and logic.] Boolean functions, truth tables;
  combinational networks and their minimisation; adders; finite-state
  machines; sequential networks; flip-flops, registers.
\item[Data representation.] Encoding; redundant codes, error detection and
  correction; decoders, multiplexers; positional systems; naturals, integers,
  reals; arithmetic.
\item[Structure of a computer.] The CPU: control unit, ALU, registers; the
  instruction cycle; the datapath; main memory; the I/O subsystem;
  microprogrammed architectures.
\item[RISC-V assembly.] Instruction set; registers, addressing; arithmetic and
  logic; branches; subroutines, calling convention, stack; memory and I/O.
\item[Project.] \vermone: a game in \riscv assembly on \prg, built in four
  guided parts, then varied, documented and packaged by the student or by a
  small team.
\end{description}

The detailed syllabus is published on the university's information system and
on the e-learning portal, \url{https://elearning.uniparthenope.it}.

\subsection{Learning outcomes}

The outcomes are stated, as usual in European higher education, along the five
Dublin descriptors.

\begin{description}[style=nextline,leftmargin=1.5em]
\item[Knowledge and understanding.] A sound knowledge of the binary system, of
  digital logic and of the operating principles of computers. No previous
  computing background is assumed.
\item[Applying knowledge.] The ability to apply theory to exercises and to
  \riscv assembly programs, and to recognise both the capabilities and the
  limits of computers on real problems.
\item[Making judgements.] The ability to write assembly programs that represent
  and process real problems in binary, choosing suitable data representations
  and control structures.
\item[Communication skills.] The ability to take an active part in teams that
  design, implement and document assembly applications, using correct
  technical terminology.
\item[Learning skills.] The ability to connect representation, logic networks
  and processors into one picture, and to keep up to date independently from
  reliable technical sources.
\end{description}

% ----------------------------------------------------------------------------
\section{Organisation and assessment}

This section records the rules of the academic year 2026/27. Dates that had
not been fixed when the book was compiled are announced in class and on the
e-learning portal, which prevail over this text in case of difference.

\subsection{Timetable}

From September to December 2026 the course meets four times a week in Aula~1:
lectures on Monday (14:00--16:00) and Wednesday (09:00--11:00), laboratory on
Tuesday and Friday (09:00--11:00). An internal calendar with the actual
distribution of lectures and laboratory sessions is shared with the students.

\subsection{One midterm test}

There is one midterm test, taken in person and on paper in about 45~minutes:
\textbf{ten multiple-choice questions plus two exercises}. It is held on a
Monday, 14:00--16:00, in Aula~1, on a date to be announced; students may be
split into groups and the test run in shifts. Candidates bring an identity
document and a pen, and must know their institutional e-mail address and
student number.

The test assesses both sides, students and teachers: it is needed to know
where the course is going. It closes the first three blocks
(Table~\ref{tab:intro:midterm}); the fourth, the machine, is examined in the
oral. \riscv instructions appear in the examples of the test: the
first-contact material of Chapter~\ref{ch:firstcontact} is assumed, and the
code of \vermone, parts 1 to 3, is a natural source of questions. Every
chapter of this book ends with questions of the same kind, and the best
preparation is to answer them week by week.

\begin{table}[htb]
\centering\small
\caption{What the midterm test covers.}\label{tab:intro:midterm}
\begin{tabularx}{\textwidth}{@{}llL@{}}
\toprule
\textbf{Block} & \textbf{Lectures} & \textbf{Topics}\\
\midrule
Information & 3 to 6 & history of computing; digital information; number representation: integers and reals\\
Logic & 8 to 11 & Boolean algebra; electronics and logic gates; combinational machines\\
State and memory & 13 to 16 & sequential machines: latches, flip-flops, state machines; registers and buses; memory and permanent storage\\
\emph{Not in the midterm} & 18 to 24 & the processor, instructions and interrupts, the video card, input and output, the outlook: part of the oral exam\\
\bottomrule
\end{tabularx}
\end{table}

\begin{prgwarn}[No tech. No GPT.]
The midterm test and the exam are \emph{\#nogpt}. No technological or human
help of any kind is allowed: no devices, no ChatGPT or other large language
models, no generative AI. Cheating cancels the test at once; more serious
behaviour has more serious consequences.
\end{prgwarn}

\subsection{Assessment}

The exam has two parts, and both must be passed.

\begin{itemize}
\item \textbf{Theory (50\%)}: an oral examination with impromptu questions on
  the whole syllabus; some answers need a diagram or a calculation on the
  board. The oral can also lower the mark obtained in the midterm test.
\item \textbf{Practical (50\%)}: the laboratory exam with Prof.\ Staffa.
\end{itemize}

The final mark is the mean of the theory mark and of the practical mark.
Midterm and exam results are kept within the current academic year.

To finalise the exam a positive mark is needed in both theory and practice,
and finalisation happens only during a theory exam session. The laboratory
exam is usually held one week before the theory session, and both sessions are
booked on \url{https://uniparthenope.esse3.cineca.it}. Two questions are asked
every year. \emph{Can the theory exam be taken before the practical one has
been passed?} Yes: the result is kept for one academic year, but the exam is
finalised only after the practical part has been passed. \emph{Can the theory
session be booked without the practical test?} Yes, the case is the same.

\subsection{Midterm week and the closing event}

During the midterm week the lectures are replaced by seminars on emerging
computer architectures. At the end of the course, a student or a team that
presents its game at the closing ceremony, a variant of \vermone or a game of
its own, can earn bonus points for the exam. The closing morning has two
halves: the last lecture, beyond one core, from superscalar \riscv to
supercomputers (Chapter~\ref{ch:beyond}), and a visit to the supercomputing
systems of the Department of Science and Technology. Finished games can be
published on the \prg Cartridge Store, as Chapter~\ref{ch:store} shows.

% ----------------------------------------------------------------------------
\section{Material}

\begin{description}[style=nextline,leftmargin=1.5em]
\item[Slides.] The slide decks that this book follows are always available in
  a shared Google Drive folder and are periodically published on
  \url{https://elearning.uniparthenope.it}. They are generated from source:
  corrections and new exercises are welcome as contributions.
\item[Platform.] \prg, free and open source, at
  \url{https://github.com/riscv-prg32/PRG32}, with its own textbook,
  \emph{Playing with RISC-V: Assembly, C, and Retro Games on the PRG32
  Platform} \citep{prg32textbook}.
\item[Textbooks.] \citet{patterson2020}, \emph{Computer Organization and
  Design, RISC-V Edition}, and \citet{tanenbaum2013}, \emph{Structured
  Computer Organization}, are the two references of the course. The syllabus
  also lists \citet{stallings2021} and \citet{bucci2017}.
\item[Standards.] The official \riscv documentation at
  \url{https://riscv.org} \citep{riscvisa}.
\end{description}

The teaching material lives in a folder shared with the students, who can
contribute to it; to obtain access a student sends the address of an account
linked to the service to the lecturer. All other communication with the
lecturer goes through the institutional address of the student.

\begin{prgaside}[Cineforum]
Thirty films and series about computing accompany the course, as a reminder
that the subject has a culture as well as a technique: \emph{The Social
Network} (2010), \emph{Steve Jobs} (2015), \emph{Pirates of Silicon Valley}
(1999), \emph{The Imitation Game} (2014), \emph{Hackers} (1995), \emph{Ex
Machina} (2014), \emph{Her} (2013), \emph{Tron} (1982), \emph{Tron: Legacy}
(2010), \emph{The Matrix} (1999), \emph{WarGames} (1983), \emph{Sneakers}
(1992), \emph{Jobs} (2013), \emph{Antitrust} (2001), \emph{Blackhat} (2015),
\emph{The Fifth Estate} (2013), \emph{The Circle} (2017), \emph{Silicon
Valley} (2014--2019), \emph{Source Code} (2011), \emph{Chappie} (2015),
\emph{Mr.\ Robot} (2015--2019), \emph{Ghost in the Shell} (1995),
\emph{Minority Report} (2002), \emph{Transcendence} (2014), \emph{Ready Player
One} (2018), \emph{A.I.\ Artificial Intelligence} (2001), \emph{Revolt}
(2017), \emph{The Thirteenth Floor} (1999), \emph{Eagle Eye} (2008) and
\emph{Terminator~2: Judgment Day} (1991). \emph{Ready Player One} is also the
cultural reference of the \prg textbook: fluency with the foundational layers
is a literacy that keeps its value.
\end{prgaside}

% ----------------------------------------------------------------------------
\section{First contact: who are you?}

The first lecture ends with questions put to the room, and a reader working
alone may put them to themselves. Their purpose is not to select but
to find where to begin.

\begin{itemize}
\item \textbf{Hardware.} Have you ever assembled a computer? Have you played
  with an Arduino or an ESP32?
\item \textbf{Software.} How many programming languages do you know? Have you
  ever read a line of assembly?
\item \textbf{Culture.} What are retrocomputing and retrogaming? Have you ever
  visited a computer museum?
\item \textbf{Concepts.} Can you tell the difference between ROM, RAM and
  storage? What is cloud computing? And the Internet of Things?
\item \textbf{RISC-V.} Have you heard of \riscv? Which processor is inside your
  telephone, your laptop, your router?
\end{itemize}

In class the questions are followed by a short anonymous questionnaire, the
\emph{Sorting Hat}, on what attracts each student most in the future of
computing and of science; its results are discussed and used to tune examples
and projects.

\section{Before the next chapter}

Chapter~\ref{ch:setup} installs everything, step by step, on Windows, macOS
and Linux. It goes faster if the long downloads are already done.

\begin{exercise}[preparation]
\begin{enumerate}[nosep]
\item Have a laptop on which you are allowed to install programs.
\item Download the \prgqt installer for your system from the \emph{Releases}
  page of \url{https://github.com/riscv-prg32/PRG32-QT}.
\item On a good connection, start the long download: ESP-IDF, following the
  \fname{README} of \url{https://github.com/riscv-prg32/PRG32} for your system.
\item If something fails, do not fight it alone: keep the error message.
\end{enumerate}
\expected{You begin Chapter~\ref{ch:setup} with the downloads done. In two
hours you have \prgqt running and a cartridge built on your own laptop.}
\end{exercise}
